\chapter{IMPLEMENTATION}\label{ch:implementation}
\begingroup\emergencystretch=3em

\section{Prototype Overview}\label{sec:prototype}
VeerDrishti is implemented as an Android application that combines a live camera feed with
on-device computer-vision inference. The implementation uses Android CameraX for video capture,
Ultralytics YOLO26 Nano and its variants for inference, Flutter for the interface, and Android
Logcat for development diagnostics. The prototype brings these components together to display
detections directly over the camera view.

The implementation focuses on the capture--inference--visualization loop. The prototype demonstrates
ordinary object detection, segmentation overlays, and human pose visualization. The longer-term
combat-vision objective guides the application, but thermal fusion, open-vocabulary recognition,
combat-specific fine-tuning, and persistent mission recording are separate extensions whose
completion is not established by the current prototype evidence.

\label{end:prototype}
\section{Module Implementation}\label{sec:modules}

\subsection{Input Module}
Android CameraX provides the live camera input. The application exposes the camera view through
the Flutter interface so that users can observe the scene and its associated overlays together.
The portrait and landscape screenshots demonstrate operation in both orientations. Camera access
requires Android camera permission; denial or interruption must be handled before attempting to
run the inference pipeline.

\subsection{Preprocessing Module}
Preprocessing is handled by the native YOLO inference pipeline. This stage includes frame rotation,
resizing or letterboxing, RGB conversion, and input tensor
preparation. Rotation keeps the analyzed image consistent with device orientation. Resizing and
letterboxing adapt the camera frame to the model input while preserving the information needed
to align output coordinates with the preview. The exact input resolution and normalization
settings depend on the deployed model configuration and should be recorded for reproducibility.

\subsection{Inference Module}
YOLO26 Nano and its task variants form the inference component. Computation takes
place on the device, supporting a prototype that does not require a remote server for each frame.
The interface exposes Det, Seg, and Pose modes, corresponding to detection, segmentation, and pose
outputs. The precise model export format, class configuration, and hardware execution delegate
must be taken from the application configuration when documenting a reproducible deployment.

\subsection{Post-Processing Module}
The native inference output is passed to the Flutter YOLOView integration for display. The
application renders bounding boxes, class labels, and confidence information on the live camera
feed. Segmentation mode adds colored object-region overlays, while pose mode displays body
keypoints and their connecting lines. Detection coordinates must be mapped to the displayed
preview after rotation, letterboxing, scaling, or zoom so that the overlays remain aligned with
the visible objects.

\subsection{User Interface Module}
Flutter supplies the application interface, camera display, mode controls, and visual output.
The demonstration screenshots show confidence and intersection-over-union (IoU) threshold controls,
zoom information, and runtime indicators. Keeping these controls in the camera view allows users
to inspect the effect of a configuration change without switching to a separate results screen.
Object labels describe model predictions; an object or person label alone does not establish
hostile intent or distinguish a combatant from a civilian.

\subsection{Development Logging Module}
Android Logcat records runtime logs and debugging information during development and testing.
These logs support diagnosis of camera initialization, inference integration, and application
behavior. This development logging mechanism is distinct from a persistent detection-event or
mission-video archive. The current implementation does not include a demonstrated database,
encrypted recording service, or post-mission retrieval interface.

\label{end:modules}
\section{Frame Processing Pipeline}\label{sec:pipeline}
The module descriptions define the following processing sequence for the prototype:
\begin{enumerate}
  \item Initialize the selected camera and model task, then present the live preview after camera
    permission is available.
  \item Obtain an analyzable frame from CameraX and account for its orientation.
  \item Resize or letterbox the frame, convert it to RGB, and prepare the input tensor in the
    native inference pipeline.
  \item Run the selected YOLO26 Nano task variant on the device and obtain its output.
  \item Apply the configured filtering settings and map the resulting coordinates to the preview.
  \item Render labels, boxes, masks, or pose keypoints through YOLOView, display runtime
    information, and continue with subsequent frames.
    \label{end:pipeline}
\end{enumerate}

\section{Interface and Detection Controls}\label{sec:controls}
The demonstration screenshots show a confidence threshold of 0.25 and an IoU threshold of 0.70.
The confidence setting controls the minimum score accepted for display. The IoU setting controls
overlap filtering where the selected model pipeline uses it. Changing either threshold can alter
the visible output, so evaluation runs should record both values along with the selected task and
model. The implementation record should state whether all task variants use the same suppression mechanism.

The interface also displays frames per second and processing times, including preprocessing,
inference, and post-processing components. These values help inspect responsiveness during a
demonstration. They are displayed runtime observations, not a substitute for repeated timing
measurements on a documented device. Chapter~\ref{ch:result} records the displayed values and
distinguishes them from a sustained-performance benchmark.

\label{end:controls}
\section{Integration and Testing}\label{sec:integration}
Integration testing connects camera capture, model execution, and overlay rendering as one
application workflow. The demonstrated exhibits involve scissors, a mouse, furniture, laptops,
bottles, and groups of people. These scenes exercise multiple visible objects and the three task
modes. Portrait and landscape exhibits also provide evidence for checking preview
orientation and output alignment.

For reproducible testing, record the Android device and version, application build, model file,
task mode, input resolution, thresholds, and execution backend. Separate functional checks from
accuracy and performance evaluation: a correctly rendered box verifies an interface path, while
precision, recall, mAP, and sustained frame rate require defined data and measurement procedures.
The available prototype material supports an initial functional demonstration; the Result chapter describes
the available observations and the evaluation details still needed.
\label{end:integration}\label{end:implementation}
\par\endgroup
